\VerdierPhysicalPageAnchor{0047}%
\VerdierSemanticAnchor{ch1-1-4-3-compatibility}
et, par suite, d'après (1.2.3.2), pour tout couple
$(\varepsilon,\sigma)$, le diagramme ci-après est commutatif :

\VerdierSemanticAnchor{ch1-1-4-3-eq8}
\begin{center}
  $\begin{tikzcd}[ampersand replacement=\&,column sep=6em,row sep=3.5em]
    FT(\sigma)
      \arrow[rr,"{m((\varepsilon,\sigma))}"]
      \arrow[dr,"{m((1,\sigma))}"']
      \& \& T'(\sigma)F \\
      \& T'(\sigma)F
        \arrow[ur,"{T'(\sigma)*m(\varepsilon)}"'] \&
  \end{tikzcd}$
  \hfill\textup{(1.4.3.8)},
\end{center}

\noindent\emph{i.e.}

\VerdierSemanticAnchor{ch1-1-4-3-eq9}
\begin{equation}
  m((\varepsilon,\sigma))=\varepsilon m((1,\sigma)).
  \tag{1.4.3.9}
\end{equation}
Posons alors :

\VerdierSemanticAnchor{ch1-1-4-3-eq10}
\begin{equation}
  m((1,\sigma))=p(\sigma)\colon
  FT(\sigma)\longrightarrow T'(\sigma)F.
  \tag{1.4.3.10}
\end{equation}

\VerdierSemanticAnchor{ch1-1-4-3-cocycle-compatibility}
La famille des isomorphismes $p(\sigma)$, $\sigma\in G$, possède alors
la propriété suivante : pour tout couple $\sigma_1,\sigma_2\in G$, le
diagramme ci-après est $\alpha_s(\sigma_1,\sigma_2)$-commutatif
(\emph{i.e.} commutatif si $\alpha_s(\sigma_1,\sigma_2)=+1$ et
anti-commutatif si $\alpha_s(\sigma_1,\sigma_2)=-1$) :

\VerdierSemanticAnchor{ch1-1-4-3-eq11}
\begin{center}
  \resizebox{0.94\textwidth}{!}{%
    $\begin{tikzcd}[ampersand replacement=\&,column sep=5em,row sep=3.5em]
      FT(\sigma_1\sigma_2)
        \arrow[rr,"{p(\sigma_1\sigma_2)}"]
        \arrow[d,"{F*c(\sigma_2,\sigma_1)}"']
        \& \& T'(\sigma_1\sigma_2)F
          \arrow[dd,"{c'(\sigma_2,\sigma_1)*F}"'] \\
      FT(\sigma_1)T(\sigma_2)
        \arrow[d,"{p(\sigma_1)*T(\sigma_2)}"']
        \& \& {} \\
      T'(\sigma_1)FT(\sigma_2)
        \arrow[rr,"{T'(\sigma_1)*p(\sigma_2)}"]
        \& \& T'(\sigma_1)T'(\sigma_2)F\mathpunct{.}
    \end{tikzcd}$%
  }
  \hfill\textup{(1.4.3.11)}
\end{center}

\VerdierSemanticAnchor{ch1-1-4-3-converse}
Réciproquement, lorsqu'on se donne un foncteur
$F\colon\mathcal C\to\mathcal C'$ et une famille d'isomorphismes
$p(\sigma)\colon FT(\sigma)\to T'(\sigma)F$, $\sigma\in G$, telle que
pour tout couple $(\sigma_1,\sigma_2)\in G\times G$ le diagramme
(1.4.3.11) soit $\alpha_s(\sigma_1,\sigma_2)$-commutatif, on peut
reconstruire d'une manière unique un $H$-foncteur :
\[
  (F,m)\colon(\mathcal C,H,T_f,c_f)
  \longrightarrow(\mathcal C',H,T'_f,c'_f)
\]
induisant l'opération canonique du groupe des signes sur $F$, en posant
(via l'identification du groupe $H$ avec le groupe décrit par
(1.4.3.6)) :

\VerdierSemanticAnchor{ch1-1-4-3-eq12}
\begin{equation}
  m((\varepsilon,\sigma))=\varepsilon p(\sigma).
  \tag{1.4.3.12}
\end{equation}

\VerdierSemanticAnchor{ch1-1-4-3-closure}
La vérification immédiate est laissée au lecteur.
